\section{Threats to Validity and Limitations}
\label{sec:limitations}

\subsection{Literature coverage}

No review of optimization can guarantee exhaustive coverage of every named method, application, or publication. The 90-entry census is a structured coverage layer, while the 42 source-backed fingerprints are representative mechanisms rather than a claim of completeness. Relevant work may be missed because of indexing, terminology, language, publication venue, or the rapid appearance of new methods. The atlas therefore treats its scope as auditable and extensible rather than closed.

Publication bias is a related concern. Positive algorithm comparisons are more likely to be emphasized than null or unfavorable results, and benchmark studies can differ substantially in tuning effort. The evidence-state model reduces but cannot eliminate this bias.

\subsection{Taxonomy subjectivity}

Mechanism normalization requires judgment. Fields such as dominant information source, memory, adaptation, diversity control, or failure signal can admit more than one defensible encoding. Structural similarity therefore depends on the selected representation and field weights. Similarity is used as an exploratory aid, not as proof of equivalence, ancestry, or lack of novelty.

The mechanism-loop abstraction is deliberately broad. Its value is comparability, but that breadth can hide domain-specific mathematical structure. Formal convergence properties and specialized problem representations must still be examined in their native theoretical frameworks.

\subsection{Evidence heterogeneity}

Published studies differ in benchmark suites, implementations, parameter tuning, stopping rules, computing environments, and statistical practice. Such heterogeneity limits quantitative pooling. The atlas therefore favors condition-specific synthesis and marks partial or contradictory evidence rather than forcing incomparable studies into a single effect estimate.

The reproducibility registry is also incomplete for some representatives. Pending fields are not interpreted as negative evidence, but they limit conclusions about the reproducibility landscape.

\subsection{Benchmark representativeness}

The project-generated experiment is restricted to noiseless continuous BBOB. BBOB provides a controlled and widely used black-box test environment, but it does not represent discrete, mixed-variable, constrained, dynamic, noisy, expensive, multiobjective, or application-specific optimization in general. Results from the project experiment should not be transferred to those domains.

Even within BBOB, a finite set of dimensions, instances, targets, and budgets samples only part of the possible condition space. Property-group interpretations depend on the validity of the benchmark-property mapping and sufficient observations within each group.

\subsection{Comparator and configuration scope}

The publication comparator panel contains four frozen configurations: Random Search, SciPy differential evolution, bounded SciPy Nelder--Mead, and pycma CMA-ES. This panel spans distinct mechanisms but is not intended to represent the full state of the art. Algorithms without an accepted publication implementation/configuration path were excluded rather than included with weaker provenance.

Performance is configuration-dependent. A result for the frozen DE or CMA-ES configuration is not a result for every variant of that family. The study also does not claim that equal objective-evaluation budgets equalize optimizer-side computational cost.

\subsection{Statistical and frontier limitations}

Discovery and held-out instance sets reduce one source of overfitting, but they do not provide unlimited external validation. More importantly, the exact reliability threshold and target subset required for a binary failure-frontier claim were not frozen before held-out execution. The study therefore does not report such a frontier. This restriction reduces claim strength but protects the distinction between predeclared confirmation and post hoc threshold selection.

No G4/G5 mechanism-repair opportunity was promoted from the frozen evidence. Consequently, the study does not test a new optimizer or repair operator. This limits claims about mechanism causality but avoids a result-driven intervention.

\subsection{Temporal validity}

Optimization research changes rapidly. The source ledger and census represent the evidence available under the review's search and freeze process, not a permanent endpoint. New algorithms, theoretical results, benchmark suites, and independent replications can change individual evidence states. The machine-readable atlas and versioned evidence model are intended to make such updates traceable.

\subsection{Reproducibility versus generalizability}

The frozen workflows strengthen computational reproducibility, but reproducibility and generalizability are distinct. A perfectly reproducible result can remain narrow if the benchmark distribution, implementation, or parameter policy is narrow. Conversely, a broad literature pattern may be difficult to reproduce exactly when original artifacts are unavailable.

The conclusions therefore remain bounded by both dimensions: what can be rerun and what the evaluated conditions can reasonably represent. This boundary informs the conclusions that follow.
